% =======================================================
% 2. TERMINI
% =======================================================
\section{Termini}
 
% ------------- A -------------
\phantomsection
\addcontentsline{toc}{subsection}{A}
\subsection*{A}
\begin{itemize}[leftmargin=*]
    \item \textbf{Amministratore:} Ruolo del gruppo responsabile di assicurare
          l'efficienza di procedure, strumenti e tecnologie a supporto del
          \textit{way of working}; redige e mantiene aggiornate le Norme di
          Progetto.
    \item \textbf{Analista:} Ruolo del gruppo responsabile delle attività di
          analisi dei requisiti; attivato dopo l'aggiudicazione del capitolato.
\end{itemize}
 
% ------------- B -------------
\phantomsection
\addcontentsline{toc}{subsection}{B}
\subsection*{B}
\begin{itemize}[leftmargin=*]
    \item \textbf{Baseline:} Versione stabile e approvata di un insieme di
          documenti o artefatti, fissata in un determinato momento del progetto
          e utilizzata come riferimento ufficiale per le fasi successive. Nel
          contesto del progetto didattico le baseline principali sono: candidatura
          (CC), Requirements and Technology Baseline (RTB) e Product Baseline (PB).
    \item \textbf{Branch:} Linea di sviluppo indipendente all'interno di un
          repository Git, utilizzata per lavorare su modifiche isolate senza
          influenzare il ramo principale (\texttt{main}).
\end{itemize}
 
% ------------- C -------------
\phantomsection
\addcontentsline{toc}{subsection}{C}
\subsection*{C}
\begin{itemize}[leftmargin=*]
    \item \textbf{Candidatura (CC):} Fase iniziale del progetto didattico, antecedente
          all'aggiudicazione del capitolato, in cui il gruppo si propone come
          fornitore producendo la Valutazione dei Capitolati e la Dichiarazione degli
          Impegni. Non è un documento ma una baseline di fase.
    \item \textbf{Ciclo Di Vita Documentale:} Sequenza di fasi che ogni documento
          percorre dalla sua creazione alla pubblicazione ufficiale. Nel contesto
          del gruppo si articola in quattro fasi: Stesura (redazione del contenuto),
          Revisione (controllo da parte del Verificatore), Pubblicazione (integrazione
          nel branch \texttt{main} tramite pull request) e Ingresso in Baseline
          (la versione diventa riferimento ufficiale).
\end{itemize}
 
% ------------- D -------------
\phantomsection
\addcontentsline{toc}{subsection}{D}
\subsection*{D}
\begin{itemize}[leftmargin=*]
    \item \textbf{Dichiarazione Degli Impegni:} Documento gestionale esterno redatto
          in fase di candidatura che specifica: scadenza ultima di consegna prevista,
          preventivo dei costi finali calcolato secondo il regolamento, totale delle
          ore produttive assegnate a persona (nell'intervallo 80--95 ore), distribuzione
          delle ore per ruolo, rischi attesi e impegni formali assunti verso il
          committente.
    \item \textbf{Discord:} Piattaforma di comunicazione vocale e testuale utilizzata
          dal gruppo come ambiente principale per le riunioni interne e il coordinamento
          quotidiano.
\end{itemize}
 
% ------------- E -------------
\phantomsection
\addcontentsline{toc}{subsection}{E}
\subsection*{E}
\textit{Nessun termine.}
 
% ------------- F -------------
\phantomsection
\addcontentsline{toc}{subsection}{F}
\subsection*{F}
\textit{Nessun termine.}
 
% ------------- G -------------
\phantomsection
\addcontentsline{toc}{subsection}{G}
\subsection*{G}
\begin{itemize}[leftmargin=*]
    \item \textbf{GitHub Pages:} Servizio di GitHub che consente di pubblicare siti
          web statici direttamente da un repository. Utilizzato dal gruppo per
          esporre la documentazione prodotta come sito web accessibile pubblicamente.
    \item \textbf{GitHub Project:} Strumento di gestione dei task integrato in
          GitHub, utilizzato per pianificare, tracciare e monitorare le attività
          del progetto tramite board e issue.
    \item \textbf{Google Calendar:} Servizio di calendario condiviso di Google,
          utilizzato dal gruppo per la gestione delle disponibilità dei membri e
          la pianificazione delle riunioni ricorrenti e straordinarie.
    \item \textbf{Google Form:} Strumento di Google utilizzato per la creazione
          guidata dei task; genera automaticamente una issue su GitHub e registra
          i dati nel foglio di calcolo condiviso per il monitoraggio delle
          misurazioni.
\end{itemize}
 
% ------------- H -------------
\phantomsection
\addcontentsline{toc}{subsection}{H}
\subsection*{H}
\textit{Nessun termine.}
 
% ------------- I -------------
\phantomsection
\addcontentsline{toc}{subsection}{I}
\subsection*{I}
\begin{itemize}[leftmargin=*]
    \item \textbf{ISO/IEC 12207:1995:} Standard internazionale
          (\textit{Information Technology~-- Software Life Cycle Processes})
          che definisce un framework di processi per il ciclo di vita del software,
          suddivisi in processi primari (es.\ fornitura, sviluppo), di supporto
          (es.\ documentazione, verifica, gestione della configurazione) e
          organizzativi (es.\ gestione, infrastruttura). Fornisce un quadro
          metodologico per strutturare, gestire e controllare le attività di
          sviluppo e manutenzione del software, applicabile tramite tailoring al
          contesto specifico del progetto.
    \item \textbf{Issue:} Unità di lavoro o segnalazione registrata su GitHub,
          utilizzata per tracciare attività, problemi o richieste. Nel flusso del
          gruppo ogni issue corrisponde a un task e viene generata automaticamente
          tramite Google Form.
\end{itemize}
 
% ------------- J -------------
\phantomsection
\addcontentsline{toc}{subsection}{J}
\subsection*{J}
\begin{itemize}[leftmargin=*]
    \item \textbf{Just-in-time:} Principio organizzativo secondo cui ogni attività
          viene normata immediatamente prima di essere attuata, evitando
          pianificazioni eccessive in anticipo. Applicato al \textit{way of working},
          implica che le norme vengano estese e aggiornate in modo incrementale,
          a intervalli brevi, man mano che nuove attività vengono introdotte.
\end{itemize}
 
% ------------- K -------------
\phantomsection
\addcontentsline{toc}{subsection}{K}
\subsection*{K}
\textit{Nessun termine.}
 
% ------------- L -------------
\phantomsection
\addcontentsline{toc}{subsection}{L}
\subsection*{L}
\begin{itemize}[leftmargin=*]
    \item \textbf{LLM (Large Language Model):} Modello avanzato di intelligenza
          artificiale addestrato su grandi quantità di testo per la comprensione,
          elaborazione e generazione di linguaggio umano in modo contestuale.
\end{itemize}
 
% ------------- M -------------
\phantomsection
\addcontentsline{toc}{subsection}{M}
\subsection*{M}
\begin{itemize}[leftmargin=*]
    \item \textbf{Misurazione:} Raccolta strutturata di dati quantitativi relativi
          allo svolgimento delle attività del gruppo, comprensiva di: ruolo
          ricoperto, tempo previsto e tempo effettivo impiegato. I dati sono
          registrati nel foglio di calcolo condiviso e utilizzati per il
          monitoraggio dell'avanzamento.
    \item \textbf{MVP (Minimum Viable Product):} Versione del prodotto dotata delle
          funzionalità minime sufficienti per essere valutata dal proponente.
          Consente di verificare la bontà della visione iniziale, di raccogliere
          feedback fondato e di prendere decisioni informate per il completamento
          definitivo del prodotto.
\end{itemize}
 
% ------------- N -------------
\phantomsection
\addcontentsline{toc}{subsection}{N}
\subsection*{N}
\begin{itemize}[leftmargin=*]
    \item \textbf{Norme Di Progetto:} Documento gestionale interno che definisce
          le procedure operative, le convenzioni redazionali e gli strumenti adottati
          dal gruppo, organizzati secondo i processi del ciclo di vita definiti da
          ISO/IEC 12207:1995. Redatto e mantenuto dall'Amministratore; soggetto
          ad aggiornamento incrementale secondo il principio just-in-time.
\end{itemize}
 
% ------------- O -------------
\phantomsection
\addcontentsline{toc}{subsection}{O}
\subsection*{O}
\begin{itemize}[leftmargin=*]
    \item \textbf{OdG Preliminare:} Acronimo di ``Ordine del Giorno Preliminare''.
          Indica l'elenco dei punti da trattare nella riunione, pubblicato dal
          Responsabile prima dell'inizio della seduta su un documento Google
          condiviso. Funge anche da base per la successiva redazione del verbale
          interno, in quanto nel medesimo documento vengono annotate le decisioni
          adottate durante la riunione.
\end{itemize}
 
% ------------- P -------------
\phantomsection
\addcontentsline{toc}{subsection}{P}
\subsection*{P}
\begin{itemize}[leftmargin=*]
    \item \textbf{Pedice G:} Convenzione tipografica adottata nelle norme di progetto
          per segnalare che un termine è definito nel Glossario. Viene applicato
          alla prima occorrenza del termine in ciascun documento tramite lo script
          Python di automazione (\texttt{python\_glossary\_automation/}).
          \textit{Esempio}: Termine\textsubscript{G}.
    \item \textbf{PoC (Proof of Concept):} Artefatto eseguibile, da considerarsi
          usa-e-getta sul piano architetturale, realizzato all'inizio del progetto
          per valutare la fattibilità tecnologica del prodotto atteso rispetto a
          specifici sottoinsiemi di funzionalità individuati con il proponente.
          Costituisce la technology baseline e aiuta a consolidare l'analisi dei
          requisiti dal lato della soluzione.
    \item \textbf{Postgres (PostgreSQL):} Sistema avanzato di gestione di database
          relazionali e ad oggetti, open source. Utilizzato come tecnologia di
          persistenza dei dati nel capitolato di riferimento.
    \item \textbf{Preventivo Dei Costi:} Stima dell'impegno economico complessivo
          del progetto, calcolata moltiplicando le ore previste per ciascun ruolo
          per il relativo costo orario ufficiale. Presentato nella Dichiarazione
          degli Impegni; il valore accettato all'aggiudicazione diventa limite
          superiore invalicabile durante lo svolgimento del progetto.
    \item \textbf{Progettista:} Ruolo del gruppo responsabile delle attività di
          progettazione architetturale (\textit{design}); attivato dopo
          l'aggiudicazione del capitolato.
    \item \textbf{Programmatore:} Ruolo del gruppo responsabile delle attività di
          codifica del prodotto software; attivato durante la fase di sviluppo.
    \item \textbf{Prompt:} Input testuale in linguaggio naturale fornito da un
          utente a un sistema di intelligenza artificiale per guidarne la
          generazione di contenuto.
    \item \textbf{Pull Request:} Meccanismo di GitHub utilizzato per proporre
          l'integrazione di modifiche da un branch di lavoro al branch
          \texttt{main}. Nel flusso documentale del gruppo costituisce il punto
          di ingresso obbligatorio per la verifica: il Verificatore esamina le
          modifiche, lascia eventuali commenti e approva o rifiuta l'integrazione.
\end{itemize}
 
% ------------- Q -------------
\phantomsection
\addcontentsline{toc}{subsection}{Q}
\subsection*{Q}
\textit{Nessun termine.}
 
% ------------- R -------------
\phantomsection
\addcontentsline{toc}{subsection}{R}
\subsection*{R}
\begin{itemize}[leftmargin=*]
    \item \textbf{Redattore:} Ruolo del gruppo responsabile della produzione del
          contenuto di un documento, nel rispetto del template approvato e delle
          convenzioni redazionali definite nelle Norme di Progetto. Non può
          verificare il proprio lavoro.
    \item \textbf{Repository:} Archivio Git centralizzato che contiene documenti,
          codice sorgente e cronologia completa delle versioni del progetto.
          Ospitato su GitHub; il branch \texttt{main} contiene esclusivamente le
          versioni approvate e stabili.
    \item \textbf{Responsabile:} Ruolo del gruppo responsabile del coordinamento
          delle attività, dell'elaborazione di piani e scadenze, dell'approvazione
          del rilascio di prodotti parziali o finali e della rappresentanza del
          gruppo verso committente e proponente.
    \item \textbf{Revisione:} Fase del ciclo di vita documentale in cui il
          Verificatore controlla struttura, coerenza, correttezza linguistica e
          conformità del documento alle Norme di Progetto, prima che esso venga
          pubblicato nel branch \texttt{main}.
    \item \textbf{RTB (Requirements and Technology Baseline):} Prima revisione
          formale di avanzamento del progetto. Fissa i requisiti da soddisfare
          (Analisi dei Requisiti) e le tecnologie adottate, dimostrate tramite
          Proof of Concept. Si articola in due parti: valutazione del docente
          Cardin e presentazione al docente Vardanega.
    \item \textbf{Ruolo:} Funzione formalmente assegnata a un membro del gruppo
          per lo svolgimento di specifiche attività. I ruoli previsti sono:
          Responsabile, Amministratore, Analista, Progettista, Programmatore,
          Verificatore. Ogni membro deve ricoprire un solo ruolo alla volta e
          ruotare su tutti i ruoli nel corso del progetto.
\end{itemize}
 
% ------------- S -------------
\phantomsection
\addcontentsline{toc}{subsection}{S}
\subsection*{S}
\begin{itemize}[leftmargin=*]
    \item \textbf{Semantic Versioning:} Convenzione di versionamento adottata dal
          gruppo basata sul formato \texttt{vX.Y.Z}, dove \texttt{X} (Major)
          identifica rilasci ufficiali o cambiamenti strutturali profondi,
          \texttt{Y} (Minor) indica aggiunte o modifiche significative di contenuto
          e \texttt{Z} (Patch) segnala correzioni ortografiche o modifiche non
          semantiche.
\end{itemize}
 
% ------------- T -------------
\phantomsection
\addcontentsline{toc}{subsection}{T}
\subsection*{T}
\begin{itemize}[leftmargin=*]
    \item \textbf{Tailoring:} Processo di adattamento di uno standard (nel caso del
          gruppo: ISO/IEC 12207:1995) al contesto specifico del progetto, mediante
          la selezione motivata dei processi da attivare e l'esclusione giustificata
          di quelli non applicabili.
    \item \textbf{Task:} Unità operativa di lavoro, identificata dal codice
          \texttt{T-y.z}, assegnata a un membro del gruppo con ruolo, scadenza,
          stima del tempo previsto e tracciamento del tempo effettivo. Creata
          tramite Google Form e tracciata su GitHub Project.
    \item \textbf{Template:} Struttura predefinita e condivisa utilizzata per la
          redazione di tutti i documenti del gruppo, al fine di garantire uniformità
          formale e coerenza tra i diversi elaborati.
    \item \textbf{Tempo Effettivo:} Tempo realmente impiegato da un membro del
          gruppo per completare un task, espresso in minuti. Registrato nel foglio
          di calcolo condiviso a consuntivo dell'attività.
    \item \textbf{Tempo Previsto:} Stima iniziale del tempo necessario per
          completare un task, espressa in minuti. Definita al momento della
          creazione del task e utilizzata come riferimento per il monitoraggio
          dell'avanzamento.
    \item \textbf{Title Case:} Convenzione tipografica adottata dal gruppo per i
          titoli di sezioni e sotto-sezioni, che prevede l'iniziale maiuscola per
          ogni parola. \textit{Esempio}: ``Ciclo Di Vita Documentale''.
    \item \textbf{T-y.z:} Codice identificativo univoco di un task, dove \texttt{y}
          è il numero progressivo del verbale in cui il task è stato originato e
          \texttt{z} è il numero progressivo del task all'interno di quel verbale.
    \item \textbf{Trunk-based Development:} Strategia di branching in cui ogni
          membro lavora su un branch dedicato a un singolo documento o attività e
          integra le modifiche nel branch principale (\texttt{main}) tramite pull
          request al termine del lavoro, eliminando poi il branch.
\end{itemize}
 
% ------------- U -------------
\phantomsection
\addcontentsline{toc}{subsection}{U}
\subsection*{U}
\textit{Nessun termine.}
 
% ------------- V -------------
\phantomsection
\addcontentsline{toc}{subsection}{V}
\subsection*{V}
\begin{itemize}[leftmargin=*]
    \item \textbf{Valutazione Dei Capitolati:} Documento gestionale esterno prodotto
          in fase di candidatura che analizza e confronta i capitolati proposti
          dalle aziende proponenti, includendo il resoconto degli incontri
          esplorativi, i criteri di valutazione, i vantaggi e gli svantaggi di
          ciascuna proposta e la motivazione della scelta finale.
    \item \textbf{VE-y.x:} Codice identificativo univoco di una decisione adottata
          in un incontro esterno (con proponente o docente), registrata nel relativo
          verbale esterno. \texttt{y} è il numero progressivo del verbale esterno
          e \texttt{x} il numero progressivo della decisione al suo interno.
    \item \textbf{Verificatore:} Ruolo del gruppo responsabile del controllo di
          conformità di documenti e codice prodotti dagli altri membri. Deve essere
          sempre un membro diverso dal Redattore dell'elaborato in esame; non può
          verificare il proprio lavoro.
    \item \textbf{VI-y.x:} Codice identificativo univoco di una decisione adottata
          in una riunione interna, registrata nel relativo verbale interno.
          \texttt{y} è il numero progressivo del verbale interno e \texttt{x} il
          numero progressivo della decisione al suo interno.
\end{itemize}
 
% ------------- W -------------
\phantomsection
\addcontentsline{toc}{subsection}{W}
\subsection*{W}
\begin{itemize}[leftmargin=*]
    \item \textbf{Way of Working (WoW):} Insieme strutturato e documentato di
          regole, strumenti, procedure e metodologie adottate dal gruppo per
          organizzare e svolgere il lavoro in modo uniforme e tracciabile.
          Corrisponde operativamente alle Norme di Progetto e si evolve in modo
          incrementale secondo il principio just-in-time per tutta la durata del
          progetto.
    \item \textbf{WhatsApp:} Applicazione di messaggistica istantanea utilizzata
          dal gruppo come canale secondario per comunicazioni rapide, notifiche
          urgenti e convocazioni di riunioni straordinarie.
\end{itemize}
 
% ------------- X -------------
\phantomsection
\addcontentsline{toc}{subsection}{X}
\subsection*{X}
\textit{Nessun termine.}
 
% ------------- Y -------------
\phantomsection
\addcontentsline{toc}{subsection}{Y}
\subsection*{Y}
\textit{Nessun termine.}
 
% ------------- Z -------------
\phantomsection
\addcontentsline{toc}{subsection}{Z}
\subsection*{Z}
\textit{Nessun termine.}